\chapter{Аналитическая часть}

В данной лабораторной работе рассматриваются два алгоритма поиска в массиве: алгоритм \textbf{поиска перебором} и алгоритм \textbf{бинарного поиска}.

\section{Алгоритм поиска перебором}

\textbf{Линейный поиск (полным перебором)} — алгоритм, который последовательно просматривает элементы массива. Поиск прекращается, если искомый элемент найден по \(i\)-му индексу или просмотрен весь массив, и элемент не найден.

Линейный поиск универсален, так как не накладывает требований к массиву, но имеет сложность \(O(N)\), что делает его не самым эффективным решением. \cite{knuth}

\section{Алгоритм бинарного поиска}

\textbf{Бинарный поиск} — алгоритм поиска в отсортированном массиве, который последовательно делит массив пополам для нахождения искомого элемента. Поиск прекращается при выполнении одного из двух условий: искомый элемент найден по \(i\)-му индексу или диапазон поиска пуст (нижний индекс превышает верхний).

Бинарный поиск требует, чтобы массив был отсортирован, и имеет сложность \(O(\log N)\), что делает его более эффективным по сравнению с линейным поиском. \cite{knuth}

\section{Вывод}

В результате аналитического раздела были рассмотрены алгоритмы линейного поиска и бинарного поиска в массивах.

\clearpage